\documentclass[11pt]{article}
\usepackage[margin=1in]{geometry}
\usepackage{amsmath,amssymb,amsthm,booktabs,enumitem,hyperref,natbib}
\hypersetup{
 colorlinks=true,
 linkcolor=blue,
 citecolor=blue,
 urlcolor=blue,
 pdftitle={Balanced transport at the egalitarian core},
 pdfauthor={Maxwell Flanagan}
}

\newtheorem{theorem}{Theorem}
\newtheorem{lemma}[theorem]{Lemma}
\newtheorem{proposition}[theorem]{Proposition}
\newtheorem{corollary}[theorem]{Corollary}
\newtheorem{remark}[theorem]{Remark}
\newtheorem{definition}[theorem]{Definition}
\newtheorem{example}[theorem]{Example}
\newcommand{\core}{\operatorname{Core}}
\newcommand{\Lmin}{\operatorname{Lmin}}
\newcommand{\Lmax}{\operatorname{Lmax}}
\newcommand{\one}{\mathbf 1}

\title{Balanced transport at the egalitarian core:\\
a complete boundary phase diagram}
\author{Maxwell Flanagan}
\date{Draft of August 4, 2026}

\begin{document}
\maketitle

\begin{abstract}
For a transferable-utility game, exact partition at an allocation means that
every cumulative union of its richest payoff strata is tight.  Both this
condition and its Lorenz consequence are due to earlier work.  We determine
exactly when the
condition can fail at the leximin core allocation of a singleton-normalized
exact game.  If a strict payoff boundary has $m$ cumulative poorer players
and $h$ richer players, then a slack richer suffix exists if and only if
\[
 m\ge m_{\min}(h),\qquad
 m_{\min}(1)=4,\quad m_{\min}(2)=m_{\min}(3)=3,
 \quad m_{\min}(h)=2\ (h\ge4).
\]
The lower bounds are structural.  A normal-cone theorem turns leximin
optimality into one simultaneous strict cover by tight coalitions.  Projecting
the cover to the richer side and applying a convex defect functional shows
more generally that every failure selects one lifted minimal balanced packet
with a positive rooted semi-balanced obstruction; mixtures are unnecessary.
The resulting certificate uses at most $|N|-1$ positive coalitions, and its
slack budget forces some one-player extension of the richer suffix to be
slack as well.
One poorer player cannot support the required opposing signs.  A unique
richest player needs four poorer players by a fractional sacrifice budget.
With two poorer players, every failure needs a proper active hyperedge of size
at least three, hence at least four richer players.  Three explicit
dimension-free lower-envelope families attain every remaining pair.  The
familiar four-player guarantee and sharp five-player failure are immediate
corollaries.  The argument enumerates no games and no payoff-stratum
patterns; its only case analyses are over the shapes of a protector trace and
over the components of a rank-two vertex.
\end{abstract}

\section{Introduction and prior-art boundary}

Let $(N,v)$ be a transferable-utility game and
\[
 \core(v)=\{x\in\mathbb R^N:x(N)=v(N),\ x(S)\ge v(S)
                   \text{ for every }S\subseteq N\}.
\]
The game is \emph{exact} if every coalition inequality is attained by some
core allocation.  Write $\phi=\Lmin(v)$ for the unique core allocation whose
increasing rearrangement is lexicographically maximal.  At a strict payoff
boundary of $\phi$, let $L$ be the cumulative poorer side and let
$H=N\setminus L$ be the richer suffix.

The exact-partition condition and its implications for Dutta--Ray and Lorenz
selection are prior art \citep{LlerenaMauri2017,
SanchezSorianoBranzeiLlorcaTijs2010}.  Qualitative tight-coalition
descriptions of the egalitarian selections are likewise established:
\citet{ArinInarra2001} characterize the egalitarian core by the impossibility
of a rich-to-poor transfer inside the core, which is exactly the existence of
a tight coalition containing the richer player and omitting the poorer one,
and they place $\Lmin$ in that set; \citet{ArinKuipersVermeulen2003}
characterize $\Lmax$ on large-core games; see also \citet{Mauri2018}.  Exact
balancedness and semi-balanced facet systems are also established
\citep{CsokaHeringsKoczy2011,LohmannBormHerings2012,
StudenyKratochvil2022}.  In particular, the latter authors already obtain
purely minimal semi-balanced systems from minimal balanced systems by
complementing the exceptional coalition.  Our candidate contribution is not
that full-set correspondence, but the Lmin-selected normal cover and its
projection to one minimal balanced packet on the richer suffix, together with
the sharp two-parameter classification and realizing families.  The closest
published low-dimensional result known to us is that exact games with at most
four players have large cores \citep{Biswas1999} and that large-core games
are egalitarian stable, with $\mathrm{PES}=\Lmax$ on that domain
\citep{Dietzenbacher2021}; that chain does not itself assert exact partition
at $\Lmin$.

The main theorem is stated now.  All games below are
singleton-normalized, meaning $v(\{i\})=0$ for every player.

\begin{theorem}[boundary phase diagram]\label{thm:phase}
At a strict payoff boundary of $\phi=\Lmin(v)$ in a singleton-normalized
exact game, put $m=|L|$ and $h=|H|$.  It is possible that
$v(H)<\phi(H)$ exactly for
\[
 m\ge m_{\min}(h),\qquad
 m_{\min}(h)=
 \begin{cases}
 4,&h=1,\\
 3,&h=2\text{ or }3,\\
 2,&h\ge4.
 \end{cases}                                      \tag{1}
\]
Every admitted pair is realized by an integral indecomposable exact game.
The poorer side may contain several payoff strata.
\end{theorem}

The proof separates a general lower-bound mechanism from three universal
constructions.  This makes the four/five threshold a corollary rather than
the theorem's organizing principle.

\section{Strict covers and balanced-trace transport}

For $\phi\in\core(v)$ write
\[
 s_\phi(S)=\phi(S)-v(S).
\]
A coalition is $\phi$-tight when its slack is zero.

\begin{lemma}[strict richer cover]\label{lem:cover}
Let $\phi=\Lmin(v)$ and let $N=L\mathbin{\dot\cup}H$ be a strict payoff
boundary.  There are nonnegative weights on proper $\phi$-tight coalitions
whose incidence load is a common positive number on $H$ and is strictly
smaller at every player in $L$.
\end{lemma}

\begin{proof}
Let $\mathcal D$ be the tangent cone of the core at $\phi$.  There is no
$d\in\mathcal D$ with $d_i\ge0$ for every $i\in L$ and one of these
inequalities strict: a sufficiently short segment $\phi+td$ remains in the
core, preserves the strict gap between $L$ and $H$, and improves the first
poorer payoff stratum on which $d$ is nonzero.

Let $K\subseteq\mathbb R^L$ be the image of $\mathcal D$ under the coordinate
projection.  It is a polyhedral cone, being the linear image of a polyhedral
cone, and the previous paragraph says $K\cap\mathbb R_+^L=\{0\}$.

We claim the polar cone
$K^\circ=\{w\in\mathbb R^L:\langle w,y\rangle\le0\text{ for all }y\in K\}$
contains a \emph{strictly} positive vector.  Ordinary separation does not
give this: two cones meeting only at the origin admit a separating $w\ne0$
with $w\ge0$, and such a $w$ may vanish on some coordinates.  Take polars
instead.  Polarity turns the intersection into a sum, so
\[
 \mathbb R^L=\{0\}^\circ=(K\cap\mathbb R_+^L)^\circ
  =\operatorname{cl}\bigl(K^\circ+(\mathbb R_+^L)^\circ\bigr)
  =\operatorname{cl}\bigl(K^\circ+\mathbb R_-^L\bigr).
\]
Both $K^\circ$ and $\mathbb R_-^L$ are polyhedral, and a sum of finitely many
polyhedral cones is polyhedral, hence closed; the closure operation is
therefore vacuous and $K^\circ+\mathbb R_-^L=\mathbb R^L$.  Write
$\one_L=w+p$ with $w\in K^\circ$ and $p\le0$.  Then $w=\one_L-p\ge\one_L>0$.
Thus $w_i>0$ for every $i\in L$ and
$\sum_{i\in L}w_id_i\le0$ for all $d\in\mathcal D$.

Extend $w$ by zero on $H$; the displayed inequality is unchanged, so $w$ lies
in the polar of $\mathcal D$ in $\mathbb R^N$.  Since
$\mathcal D=\{d:d(N)=0,\ d(T)\ge0\text{ for every proper }\phi\text{-tight }T\}$,
that polar is generated by $\pm\one_N$ and by $-\one_T$, giving
\[
 w=\mu\one_N-\sum_T\lambda_T\one_T,
 \qquad \lambda_T\ge0,                              \tag{NC}
\]
where $T$ ranges over proper $\phi$-tight coalitions.  Thus the weighted
incidence load is $\mu$ on $H$ and $\mu-w_i<\mu$ at $i\in L$.  Nonnegativity
of the load and positivity of $w_L$ imply $\mu>0$.  Deleting coalitions
disjoint from $H$ preserves the richer load and only lowers poorer loads.
\end{proof}

Normalize the common rich load in Lemma~\ref{lem:cover} to one and remove
coalitions disjoint from $H$.  Project every remaining tight coalition to
its nonempty trace on $H$.  The projected weights lie in the exact-cover
polytope
\[
 \mathcal P_H=\left\{\beta\ge0:
   \sum_{\emptyset\ne R\subseteq H}\beta_R\one_R=\one_H\right\}.
                                                               \tag{2}
\]
Its vertices are the unique normalized weights of minimal balanced
collections on $H$; this is standard balanced-collection geometry
\citep{LohmannBormHerings2012,LaplaceMermoudEtAl2025}.

\begin{theorem}[balanced-trace transport]\label{thm:transport}
Suppose $s=s_\phi(H)>0$.  Call a \emph{lifted packet} a pair consisting of a
vertex $\beta$ of $\mathcal P_H$ together with a choice, for each trace $R$
in its support, of one supported $\phi$-tight coalition whose trace on $H$ is
$R$; its \emph{poor-load vector} $r\in\mathbb R_+^L$ records the total weight
of chosen coalitions containing each $i\in L$.  Then there are a core point
$z$, a displacement $\Delta=z-\phi$, and a finite convex mixture
$\{(\theta_q,r^q)\}$ of lifted packets such that
\begin{align}
 \Delta(H)&=-s,&\Delta(L)&=s,                                      \tag{3a}\\
 \min_{i\in L}\Delta_i&<0<\max_{i\in L}\Delta_i,                  \tag{3b}\\
 -s+\sum_{i\in L}r_i^q\Delta_i&\ge0&&\text{for every }q,           \tag{3c}\\
 \sum_q\theta_qr_i^q&<1&&\text{for every }i\in L.                 \tag{3d}
\end{align}
\end{theorem}

\begin{proof}
Exactness supplies $z\in\core(v)$ with $z(H)=v(H)$, proving (3a).
Let $c_i<1$ be the poorer loads of a normalized cover from
Lemma~\ref{lem:cover}.  Every supported coalition is tight at $\phi$ and
feasible at $z$, hence
\[
 0\le\sum_T\lambda_T\Delta(T)
   =-s+\sum_{i\in L}c_i\Delta_i.                    \tag{4}
\]
If every poorer displacement were nonnegative, the right side would be
strictly smaller than $-s+\Delta(L)=0$.  Since $\Delta(L)=s>0$, both signs
occur, proving (3b).

Let $\beta_R=\sum_{T\cap H=R}\lambda_T$ be the projected weights, so
$\beta\in\mathcal P_H$, and decompose $\beta=\sum_p\theta_p\beta^p$ into
vertices.  Independently for each trace $R$, choose a preimage $T$ with
$T\cap H=R$ with conditional probability $\lambda_T/\beta_R$.  A vertex
together with such a section is a lifted packet, and the product measure
assigns it the weight $\theta_p\prod_R(\lambda_{T_R}/\beta_R)$; index the
resulting finite mixture by $q$.

Fix one lifted packet $q$, with vertex $\beta^p$ and section $\sigma$.  Each
$\sigma(R)$ is tight at $\phi$ and feasible at $z$, so
$\sum_R\beta^p_R\Delta(\sigma(R))\ge0$.  A rich player $j$ lies in
$\sigma(R)$ exactly when $j\in R$, so its load is
$\sum_{R\ni j}\beta^p_R=1$, while a poor player $i$ has load $r_i^q$ by
definition.  The displayed sum therefore equals
$\Delta(H)+\sum_{i\in L}r_i^q\Delta_i=-s+\sum_{i\in L}r_i^q\Delta_i$,
proving (3c) for every packet, not merely on average.

For (3d), average $r_i$ over the product measure.  Since
$\sum_p\theta_p\beta^p_R=\beta_R$ and the section is chosen independently per
trace,
\[
 \sum_q\theta_qr_i^q
 =\sum_R\beta_R\cdot
   \frac{1}{\beta_R}\!\!\sum_{T\cap H=R,\ i\in T}\!\!\lambda_T
 =\sum_{T\ni i}\lambda_T=c_i<1 .
\]
\end{proof}

The mixture in Theorem~\ref{thm:transport} always contains one packet that
already carries a positive rooted obstruction.  This gives a general
dimension-free reduction before the cardinality arguments begin.

\begin{theorem}[one-packet reduction]\label{thm:packet}
At every slack strict boundary there is a lifted vertex of
$\mathcal P_H$, with poorer load vector $r\ge0$, such that, on putting
\[
 b=\max_{i\in L}r_i,\qquad u_i=b-r_i,
 \qquad \delta=1+\sum_{i\in L}u_i-b,
\]
one has $\delta>0$ and
\[
 \sum_T\lambda_T\one_T
 +\sum_{i\in L}u_i\one_{H\cup\{i\}}
 -\delta\one_H=b\one_N.                            \tag{P}
\]
Here the $T$ are proper $\phi$-tight coalitions and their traces form one
minimal balanced collection on $H$.  Moreover,
\[
 \delta s_\phi(H)
 \le\sum_{i\in L}u_i s_\phi(H\cup\{i\}),           \tag{PB}
\]
so at least one $u_i$ is positive.  The trace collection is either the
full-trace packet $\{H\}$ or consists entirely of proper subsets of $H$.
The identity uses at most $|H|$ packet coalitions and at most $|L|-1$
spokes, hence at most $|N|-1$ positive coalitions.
\end{theorem}

\begin{proof}
Let $c_i=\sum_q\theta_qr_i^q<1$ be the average loads in (3d).  For a
nonnegative vector $a$ define
\[
 F(a)=1+(|L|-1)\max_i a_i-\sum_i a_i.
\]
This function is convex, and
$F(c)\ge1-\max_i c_i>0$.  Jensen's inequality therefore gives a lifted
vertex with $F(r)>0$; this is exactly the displayed $\delta$.  Its rich load
is one.  In (P), poor player $i$ has load $r_i+u_i=b$, while every rich
player has load $1+\sum_i u_i-\delta=b$.

The positive incidence columns at a vertex of $\mathcal P_H$ are linearly
independent: otherwise a sufficiently small signed kernel perturbation would
remain nonnegative in both directions.  Thus the packet has at most $|H|$
traces.  The pure lift uses one coalition per trace.  Any $i$ attaining the
maximum $b$ has $u_i=0$, so at most $|L|-1$ spokes are positive, proving the
sparsity bound.

Choose $z\in\core(v)$ exact on $H$.  Evaluate (P) at $z$ and at $\phi$.
Writing $\Delta=z-\phi$, the packet terms satisfy $\Delta(T)\ge0$, the
spokes satisfy
$\Delta(H\cup\{i\})\ge-s_\phi(H\cup\{i\})$, and
$\Delta(H)=-s_\phi(H)$.  The displacement form of (P) is therefore
\[
 0=\sum_T\lambda_T\Delta(T)
   +\sum_i u_i\Delta(H\cup\{i\})-\delta\Delta(H)
 \ge-\sum_i u_i s_\phi(H\cup\{i\})+\delta s_\phi(H),
\]
which is (PB).  If every $u_i$ were zero, (PB) would contradict
$\delta s_\phi(H)>0$.  Finally, a minimal balanced collection containing $H$
must be the singleton collection $\{H\}$; otherwise all its traces are proper.
\end{proof}

Thus every failure has an exterior full-trace packet or an internal minimal
balanced-transport packet.  The incidence identity is an established
semi-balanced system with exceptional coalition $H$; the new point proposed
here is that the Lmin normal certificate and the convex functional $F$
select one minimal balanced packet on $H$ from an arbitrary failure.  This is
different from the established complement-of-the-exceptional-set
correspondence, which operates on the full player set.

\begin{remark}[locating the packets in the published catalogue]
\label{rem:catalogue}
Every system (P) is semi-balanced with exceptional coalition $H$ in the sense
of \citet{StudenyKratochvil2022}, and is minimal in every instance computed
for this paper.  Their electronic catalogue lists the permutational types of the
\emph{indecomposable} min-semi-balanced systems on at most six players,
sixteen on five players and fifty-three on six; these are exactly the types
that index facets of the cone of exact games.  Indecomposability is here a
property of the set system and is unrelated to direct-sum indecomposability
of a game.  Five of the five-player types and fifteen of the six-player types
have a singleton exceptional coalition, which is necessary for an $h=1$
obstruction.  It is not sufficient: at $h=1$ the exact-cover polytope (2) has
the single vertex $\{H\}$ with weight one, so a system (P) has one lift, at
most one positive coalition of size exceeding two, and all positive
coefficients equal to one.  Only three of the five singleton-exceptional
five-player types have that form; the other two are never systems (P), at any
boundary of any exact game on any number of players.

Take the smallest realization of each of the three minimal pairs of the
staircase (1).  For the $(4,1)$ member of Proposition~\ref{prop:family1} with
richest player $e$, the tight proper coalitions at $\phi$ are $ce$, $de$,
$cde$, $acde$, $bcde$, and every admissible packet coalition produces a
catalogued type:
\[
\begin{array}{ll}
 T=ce\text{ or }de: & \one_{ae}+\one_{be}+\one_{ce}+\one_{de}-3\cdot\one_e
                      =\one_N,\\
 T=cde:             & \one_{ae}+\one_{be}+\one_{cde}-2\cdot\one_e=\one_N,\\
 T=acde:            & \one_{acde}+\one_{be}-\one_e=\one_N.
\end{array}
\]
For the $(3,2)$ member of Proposition~\ref{prop:family3}, where
$\phi=(1,1,2,3,3)$ and $H=de$, the only tight proper coalitions are $ade$ and
$cde$; both lifts give the full-trace packet, and the single resulting type
is again catalogued:
\[
 \one_{ade}+\one_{bde}+\one_{cde}-2\cdot\one_{de}=\one_N.
\]
For the $(2,4)$ member of Proposition~\ref{prop:family2}, where
$\phi=(2,2,4,4,4,4)$ and $H=cdef$, three types occur:
\[
\begin{array}{ll}
 \text{catalogued:} &
   \one_{acdef}+\one_{bcdef}-\one_{cdef}=\one_N,\\
 \text{not catalogued:} &
   \one_{ac}+\one_{def}+\one_{bcdef}-\one_{cdef}=\one_N,\\
 \text{not catalogued:} &
   \one_{cef}+\one_{def}+\one_{acde}+\one_{acdf}
   +2\cdot\one_{bcdef}-3\cdot\one_{cdef}=2\cdot\one_N.
\end{array}
\]
The last two systems are min-semi-balanced but decomposable, by $E=\{c\}$ and
by $E=\{c,d,e\}$; for the second system, for instance,
$\one_{ac}+\one_{bcdef}-\one_c=\one_N$ is a combination with nonnegative
coefficients inside the system and a negative one on $E$.  Two points follow.
First, the certificate of Theorem~\ref{thm:packet} is
always a valid inequality for the exact cone but need not be a facet, so
being catalogued is strictly stronger than being a system (P).  In the
instances computed, non-facet certificates arose only in the two-poorer
family.  Second, a single failure does not select a single class: which class
appears depends on the tight coalition the cover supports.  This is the
concrete form of the multi-class phenomenon noted in Section~6.
Decomposability is a statement about set systems alone; here $bcdef$ carries
positive slack at $\phi$, so it exhibits no sharper certificate at this
boundary.
\end{remark}

Two consequences are worth isolating.  The first records what the certificate
costs; the second is a propagation statement that needs no cardinality
hypothesis.

\begin{corollary}[linear-size witness and slack propagation]
\label{cor:certificate}
Let a strict payoff boundary of $\phi=\Lmin(v)$ have $s_\phi(H)>0$.  Then:
\begin{enumerate}[nosep,label=(\roman*)]
\item the obstruction is exhibited by at most $|N|-1$ positive coalitions
      together with $H$, and by the rational data
      $\lambda,u,\delta,b$ of (P);
\item $\displaystyle
      s_\phi(H)\le\frac1\delta\sum_{i\in L}u_is_\phi(H\cup\{i\})$;
\item consequently $s_\phi(H\cup\{i\})>0$ for at least one $i\in L$.
\end{enumerate}
\end{corollary}

\begin{proof}
Part (i) is the sparsity bound of Theorem~\ref{thm:packet}.  Part (ii) is
(PB) divided by $\delta>0$.  For (iii), the right side of (ii) is positive
because the left side is, and every $u_i$ is nonnegative, so some product
$u_is_\phi(H\cup\{i\})$ is positive.
\end{proof}

Part (iii) says that a slack richer suffix is never isolated: the failure
propagates to a one-player enlargement of $H$.  Since the enlargement is
still a proper coalition, iterating is not possible in general, and the
statement is genuinely about the single step.  Part (ii) quantifies this: the
slack at $H$ is financed, at exchange rate $\delta$, by the slacks of at most
$|L|-1$ single-player completions.

The first exclusion is immediate.

\begin{corollary}\label{cor:m1}
A slack boundary cannot have $m=1$.
\end{corollary}

\begin{proof}
Equation (3b) requires both a positive and a negative displacement on $L$.
\end{proof}

\section{The two remaining lower bounds}

\subsection{A unique richest player}

Assume $H=\{r\}$, write $\phi_r=\gamma$, and put
$q_i=\phi_i<\gamma$ for $i\in L$.  The transfer-blocking characterization of
the egalitarian core \citep{ArinInarra2001} gives, for each $i\in L$, a tight
protector
\[
 T_i=\{r\}\cup Q_i,\qquad
 \emptyset\ne Q_i\subseteq L\setminus\{i\}.          \tag{5}
\]
The nonemptiness follows because the normalized singleton $\{r\}$ is slack.
Alternatively (5) is the case $H=\{r\}$ of Lemma~\ref{lem:cover}: the cover
has load $\mu$ at $r$ and load $\mu-w_i<\mu$ at $i$, so the supported
coalitions containing $r$ but not $i$ carry total weight $w_i>0$ and at least
one of them is positive.  The argument below therefore does not depend on the
cited characterization.

\begin{lemma}[sacrifice budget]\label{lem:sacrifice}
For weights $\lambda_i\ge0$ satisfying
\[
 c_j:=\sum_{i:j\in Q_i}\lambda_i\le1,\qquad
 \Lambda:=\sum_i\lambda_i,
\]
one has
\[
 \gamma(\Lambda-1)\le\sum_{j\in L}(1-c_j)q_j.        \tag{6}
\]
\end{lemma}

\begin{proof}
Exactify $\{r\}$ at a core point $z$ and put $d=z-\phi$.  Summing the
tight protector inequalities with the displayed weights, using efficiency
and the singleton lower bounds $z_j\ge0$, gives (6).  Equivalently, complete
the uncovered poorer loads by singleton coalitions; together with negative
weight $\Lambda-1$ on $\{r\}$, the incidence vector is exact-balanced, and
the exact-balanced slack inequality is exactly (6).
\end{proof}

\begin{proposition}\label{prop:h1}
If $h=1$ and the suffix is slack, then $m\ge4$.
\end{proposition}

\begin{proof}
For $m=1$ no protector in (5) exists.  For $m=2$, both protector traces are
the opposite singleton; weight both by one in (6), obtaining
$\gamma\le0$.

Let $m=3$.  If some $Q_i=\{j\}$, weight the protectors indexed by $i$ and
$j$ by one.  Their traces are disjoint and leave at most one player $u$
uncovered, so (6) gives either $\gamma\le0$ or
$\gamma\le q_u<\gamma$.  If no $Q_i$ is a singleton, then
$Q_i=L\setminus\{i\}$ for all $i$.  Weight all three protectors by $1/2$;
every poorer player has load one and (6) gives $\gamma/2\le0$.
\end{proof}

\subsection{Exactly two poorer players}

We first isolate the exact-face inequality used in the rounding argument.

\begin{lemma}[exact-face overload]\label{lem:overload}
Let $v$ be singleton-normalized and exact, let $\phi\in\core(v)$, and
let $D\subseteq H\subsetneq N$. Suppose nonnegative weights on proper
$\phi$-tight coalitions have load $1+q$ on $D$, load one on $H\setminus D$,
and loads $0\le c_i\le1$ on $i\in N\setminus H$, where $q>0$. Then
\[
 q s_\phi(D)\le\sum_{i\notin H}(1-c_i)\phi_i.       \tag{7}
\]
\end{lemma}

\begin{proof}
Choose $x\in\core(v)$ with $x(D)=v(D)$ and put $d=x-\phi$.
Summing the supported tight-coalition inequalities at $x$ gives
\[
 0\le d(H)+q d(D)+\sum_{i\notin H}c_i d_i.
\]
Efficiency replaces $d(H)$ by $-\sum_{i\notin H}d_i$. Hence
\[
 q s_\phi(D)=-q d(D)
 \le\sum_{i\notin H}(1-c_i)(-d_i)
 \le\sum_{i\notin H}(1-c_i)\phi_i,
\]
where the last inequality uses $x_i\ge v(\{i\})=0$.
\end{proof}

Put $L=\{a,b\}$ and label the displacement in
Theorem~\ref{thm:transport} so that $\Delta_a<0<\Delta_b$.  By (3d), some
lifted extreme packet has $r_b<1$.

\begin{lemma}[proper-hyperedge obstruction]\label{lem:rankthree}
Every such underloaded packet contains a proper rich trace $R$ with
\[
 3\le |R|\le |H|-1.                                  \tag{8}
\]
\end{lemma}

\begin{proof}
First dispose of the full-trace packet, for which the assertion would be
vacuous.  If $H$ is itself a supported trace then $\{H\}$ is already balanced
on $H$, so minimality forces the packet to be $\{H\}$ and there is no proper
trace at all.  This cannot happen at an underloaded packet.  Indeed the
single lifted coalition $T$ is proper with $H\subseteq T$, hence
$T\in\{H,\ H\cup\{a\},\ H\cup\{b\}\}$ and the poor load is
$r\in\{(0,0),(1,0),(0,1)\}$.  Underloading gives $r_b<1$, so $r_b=0$, and
(3c) becomes $-s+r_a\Delta_a\ge0$.  That fails for $r_a=0$ because $s>0$,
and fails for $r_a=1$ because $\Delta_a<0$.  Every supported trace of the
selected packet is therefore a proper subset of $H$.

Now suppose that every such trace has size at most two.
The positive incidence columns of a vertex of (2) are independent.  Its
rank-two components are therefore isolated singletons of weight one,
isolated edges of weight one, and vertex-disjoint odd cycles whose edges
have weight $1/2$.  This is the classical description of the vertices of a
fractional perfect matching polytope, and follows by leaf deletion.  In a
component, independence forces at most as many positive traces as vertices.
A degree-one vertex $j$ lies in a unique trace $R$, and the balance equation
at $j$ forces $\beta_R=1$; the balance equations at the other endpoint of
$R$ then kill every remaining trace there, so $R$ splits off as an isolated
edge, or as an isolated singleton when $R=\{j\}$.  Iterating removes all
leaves, so a positive singleton cannot persist inside a larger component and
every surviving component is a single cycle.  An even cycle has dependent
incidence columns, hence is excluded; an odd cycle has nonsingular incidence
matrix whose unique balanced solution is the uniform weight $1/2$.

An integral packet has integer $b$-load, and (3c) forces that load to be at
least one.  Hence the selected packet is fractional.  Its poor loads are
half-integers.  From
\[
 -s+r_a\Delta_a+r_b\Delta_b\ge0,\qquad
 \Delta_a+\Delta_b=s,
\]
and $r_b<1$, one obtains $r_b=1/2$ and $r_a=0$.  Exactly one half-weight
edge on one odd cycle therefore carries $b$, and no supported coalition
carries $a$.

Round that labeled odd cycle alternately, starting with the labeled edge,
and leave every other component at its original weight.  The resulting
nonnegative combination of proper $\phi$-tight coalitions has rich load
$\one_H+e_d$ for one rich player $d$, poor load $e_b$, and no load on $a$.
Exactify singleton $d$.  Summing the selected tight inequalities and adding
the omitted singleton $a$ gives the exact-face overload inequality
$\phi_d\le\phi_a$, contradicting the strict boundary. Indeed,
Lemma~\ref{lem:overload} has $D=\{d\}$, $c_b=1$, and $c_a=0$.
Thus (8) holds.
\end{proof}

\begin{proposition}\label{prop:m2}
If $m=2$ and the suffix is slack, then $h\ge4$.
\end{proposition}

\begin{proof}
Lemma~\ref{lem:rankthree} requires a proper trace of size at least three.
\end{proof}

The proof so far yields every forbidden pair in (1) without enumerating
games or payoff-stratum patterns.

\section{Three universal realizing families}

We use a standard lower-envelope observation.

\begin{lemma}[lower-envelope realization]\label{lem:envelope}
Let $G$ be a finite family of nonnegative allocations with common total
$T$, and define $v(S)=\min_{g\in G}g(S)$.  Then every $g\in G$ belongs to
$\core(v)$, every coalition attains its worth, and $v$ is exact.  If every
coordinate is zero in some row, then $v$ is singleton-normalized.
\end{lemma}

\begin{proof}
The defining inequality gives $g(S)\ge v(S)$ for every $g,S$, so every row
is a core allocation.  Conversely every coalition minimum is attained at a
row.  The last assertion is immediate.
\end{proof}

We record the families compactly.  The displayed face inequalities give
short progressive leximin proofs; no search is involved.

\begin{proposition}[unique-richest family]\label{prop:family1}
For every $m\ge4$ there is an integral indecomposable exact game with
$m$ poorer players $P=\{1,\ldots,m\}$, one richer player $r$, and
\[
 \Lmin(v)=\one_P+(m-2)e_r,\qquad v(\{r\})=0.
\]
\end{proposition}

\begin{proof}
Apply Lemma~\ref{lem:envelope} to the four common-total rows
\begin{align*}
 g^1&=(m-1)(e_4+e_r),\\
 g^2&=(m-1)(e_3+e_4),\\
 g^3&=e_1+e_2+e_3+(m-4)e_5+(m-1)e_r,\\
 g^4&=\one_P+(m-2)e_r,
\end{align*}
where the zero $e_5$ term is omitted at $m=4$.
Every coordinate vanishes in a row.  The four attained faces
\[
 v(N\setminus\{1\})=v(N\setminus\{2\})=2m-3,\qquad
 v(\{3,r\})=v(\{4,r\})=m-1
\]
show that any core point giving everyone at least one fixes players $1,2$
at one; subtracting $\one_N$ and comparing the two pair inequalities with
efficiency then fixes every poorer player at one.  The remaining total gives
$x_r=m-2$.  Hence $g^4=\Lmin(v)$.  Row $g^2$ proves $v(\{r\})=0$.
For indecomposability, observe that
\[
 v(\{1,4\})=v(\{2,4\})=v(\{3,4\})=1,
 \qquad v(\{3,r\})=v(\{4,r\})=m-1,
\]
and, when $m\ge5$, $v(\{4,5\})>0$.  These positive pair worths place
players $1,\ldots,5,r$ in one direct-sum component.  If $m=6$, the coalition
$T=\{3,4,5,r\}$ satisfies $v(T)=7<v(T\cup\{6\})=8$.  If $m\ge7$, put
$T=\{4,5\}$; then $v(T)=2<v(T\cup\{j\})=3$ for every $j\ge6$.
Since every remaining singleton has value zero, these strict interactions
attach every remaining player to the same component.
\end{proof}

\begin{proposition}[two-poorer internal-sacrifice family]
\label{prop:family2}
For every $m\ge2$ and $h\ge4$ there is an integral indecomposable exact game
with a poorer block $P$ of size $m$, a richer suffix $H$ of size $h$, and
\[
 \Lmin(v)=D\one_P+2D\one_H,\qquad
 v(H)=2D(h-1),\qquad D=(h-2)(h-3).
\]
\end{proposition}

\begin{proof}
Put $r=h-2$, write $P=\{p_1,p_2\}\mathbin{\dot\cup}Q$ and
$H=\{a,b\}\mathbin{\dot\cup}K$, $|K|=r$.  Use the common-total rows
\begin{align*}
 u&=D e_{p_2}+D\one_Q+3D(e_a+e_b)
       +(2r-1)(r-1)\one_K,\\
 z&=4D e_{p_1}+D\one_Q+2(r^2-1)\one_K,\\
 w^k&=D e_{p_1}+D\one_Q+3D(e_a+e_b)
       +r(2r-1)\one_{K\setminus\{k\}}\quad(k\in K),\\
 \phi&=D\one_P+2D\one_H,\\
 q^j&=\phi+D(e_{p_1}-e_j)\quad(j\in Q).
\end{align*}
Lemma~\ref{lem:envelope} gives a normalized integral exact game.  Direct
substitution yields
\begin{align*}
 v(N\setminus\{p_2\})&=D(m+2h)-D,\\
 v(H\setminus\{a\})=v(H\setminus\{b\})&=2D(h-1),\\
 v(\{p_1,a,b\}\cup(K\setminus\{k\}))&=D(2h-1),\\
 v(H)&=2D(h-1).
\end{align*}
At a core point $x$ above the floor $D$, the complement face first fixes
$x_{p_2}=D$. Write $y=x-D\one_N$, $E=y(Q)$, and $t=y(K)$. Then
$y\ge0$, $y_{p_2}=0$, and $y(N)=hD$. The two co-singleton faces give
\[
 y_a+t\ge(h-1)D,\qquad y_b+t\ge(h-1)D.             \tag{9}
\]
The $r=h-2$ remaining faces give, for every $k\in K$,
\[
 y_{p_1}+y_a+y_b+t-y_k\ge(h-1)D.
\]
Summing them and substituting
$y_{p_1}+E+y_a+y_b+t=hD$ yields $t+rE\le rD$. On the other hand,
(9) and efficiency imply
\[
 y_{p_1}+E\le t-rD\le-rE.
\]
Nonnegativity forces $y_{p_1}=E=0$. Hence all of $P$ is pinned at level
$D$; efficiency
leaves total $2Dh$ on $H$, equalized by $\phi$.  Thus $\phi=\Lmin(v)$ and
its richer suffix has slack $2D$.  The same construction is indecomposable.
Indeed, $p_1$ has positive pair worth with every other player:
\[
 v(\{p_1,p_2\})=v(\{p_1,j\})=v(\{p_1,k\})=D
 \quad(j\in Q,k\in K),
 \qquad
 v(\{p_1,a\})=v(\{p_1,b\})=3D.
\]
Thus the positive-pair graph is connected.
\end{proof}

\begin{proposition}[three-poorer escape family]\label{prop:family3}
For every $m\ge3$ and $h\ge2$ there is an integral indecomposable exact game
with $m$ cumulative poorer players and richer suffix $H$, $|H|=h$, such that
\[
 \Lmin(v)=\one_{\{p_1,p_2\}\cup Q}+2e_{p_3}+3\one_H,\qquad
 v(H)=3h-1,\qquad |Q|=m-3.
\]
\end{proposition}

\begin{proof}
Write $H=\{r_1,\ldots,r_h\}$ and take the lower envelope of
\begin{align*}
 a&=2e_{p_1}+3e_{p_3}+\one_Q+5e_{r_1}
       +3\one_{H\setminus\{r_1,r_h\}},\\
 b&=2e_{p_2}+e_{p_3}+\one_Q+4e_{r_1}
       +3\one_{H\setminus\{r_1\}},\\
 c&=\one_Q+(3h+4)e_{r_h},\\
 \phi&=e_{p_1}+e_{p_2}+2e_{p_3}+\one_Q+3\one_H,\\
 q^j&=\phi+e_{p_1}-e_j\quad(j\in Q).
\end{align*}
All rows have total $m+3h+1$, and every coordinate vanishes in a row, so
Lemma~\ref{lem:envelope} applies.  The attained values
\begin{align*}
 v(N\setminus\{j\})&=m+3h\quad(j\in Q),\\
 v(H\cup\{p_3\})&=3h+2,\\
 v(H\cup\{p_1\})&=3h+1,\\
 v(H)&=3h-1
\end{align*}
give the progressive proof.  On the face where every player receives at
least one, the complement faces fix $Q$, the second face fixes $p_1,p_2$ at
one, and the third fixes $p_3$ at two.  Efficiency leaves total $3h$ on
$H$, equalized at three by $\phi$.  Thus $\phi=\Lmin(v)$ and $s_\phi(H)=1$.
When $Q\ne\emptyset$, one member of $Q$ has positive pair worth with every
other player.  When $Q=\emptyset$, the positive-pair graph connects all but
$p_2$, and
\[
 v(\{p_1,p_3,r_h\})=4<5=v(\{p_1,p_2,p_3,r_h\})
\]
attaches $p_2$.  Hence the game is indecomposable.
\end{proof}

\begin{proof}[Proof of Theorem~\ref{thm:phase}]
Corollary~\ref{cor:m1} excludes $m=1$ for every $h$.
Proposition~\ref{prop:h1} excludes $m\le3$ when $h=1$, and
Proposition~\ref{prop:family1} realizes every $m\ge4$.
Proposition~\ref{prop:m2} excludes $h\le3$ when $m=2$, and
Proposition~\ref{prop:family2} realizes every $m\ge2,h\ge4$.
All remaining pairs have $m\ge3,h\ge2$ and are realized by
Proposition~\ref{prop:family3}.
\end{proof}

\section{The four/five threshold and its mechanism}

\begin{corollary}[sharp player threshold]\label{cor:fourfive}
Every singleton-normalized exact game on at most four players is exact
partition at $\Lmin$.  Five players suffice for failure.
\end{corollary}

\begin{proof}
Every pair $(m,h)$ with $m+h\le4$ is forbidden by (1).  The pair
$(m,h)=(3,2)$ is realized by Proposition~\ref{prop:family3} on five players.
Alternatively, $(4,1)$ is realized by Proposition~\ref{prop:family1}.
\end{proof}

The theorem gives two distinct first mechanisms.  The $(4,1)$ obstruction
is an exterior sacrifice: the fourth poorer player supplies the first
positive singleton budget capable of financing a slack richest singleton.
The $(3,2)$ obstruction is a mixed block-star mechanism.  Two-poorer
failures do not yet exist on five players: their active balanced packet must
contain a proper rank-three trace, which first fits at $(2,4)$ on six
players.  Thus ``five'' is not an unexplained output of enumeration; it is
the first total size at which either of the two admissible boundary regions
in (1) is entered.

\section{Which system the selection produces}
\label{sec:selection}

Theorem~\ref{thm:packet} produces a system (P) but says nothing about which
system.  This section answers that question completely on five players and
partially in general.  Throughout, $v$ is singleton-normalized and exact,
$\phi=\Lmin(v)$, $N=L\mathbin{\dot\cup}H$ is a strict payoff boundary with
$s=s_\phi(H)>0$, $m=|L|$ and $h=|H|$.  Two conventions are used repeatedly and
recorded once.  Core allocations of a normalized game are nonnegative, so
$\phi_i\ge0$ for every $i$ and hence $\phi_j>0$ for $j\in H$; consequently no
singleton $\{j\}$ with $j\in H$ is $\phi$-tight, and $H$ itself is not
$\phi$-tight.  In (P) the coefficient of $H$ is $-\delta<0$ and every other
coefficient is nonnegative, so $H$ is the unique exceptional set of the system
and the catalogue's constant $r$ is the quantity $b$.

It is convenient to name the objects that Theorem~\ref{thm:transport} produces
before Theorem~\ref{thm:packet} selects among them.

\begin{definition}[normal certificate]\label{sel:def:normal}
A \emph{normal certificate} at the boundary $L|H$ is a nonnegative weighting
$\lambda$ of proper $\phi$-tight coalitions, each meeting $H$, whose incidence
load equals one at every $j\in H$; writing $c_i$ for the load at $i\in L$,
$b=\max_i c_i$, $u_i=b-c_i$ and $\delta=1+\sum_iu_i-b$, one requires
$\delta>0$.  Its \emph{system} is
\[
 \Lambda=\sum_T\lambda_T\one_T+\sum_{i\in L}u_i\one_{H\cup\{i\}}-\delta\one_H
        =b\,\one_N .
\]
\end{definition}

The normalized cover of Lemma~\ref{lem:cover} is a normal certificate, because
its poorer loads satisfy $c_i<1$ and therefore
$\delta=F(c)\ge1-\max_ic_i>0$; and the packet of Theorem~\ref{thm:packet} is
the special normal certificate whose $\lambda$ is a lifted vertex of
$\mathcal P_H$.  An exclusion proved for normal certificates is thus
\emph{robust}: it survives both the mixture of Theorem~\ref{thm:transport} and
the Jensen selection.

\subsection{Three elementary lemmas}

The first is the qualitative content of the slack budget (PB).

\begin{lemma}[positive load]\label{sel:lem:spoke}
Every normal certificate at a slack boundary has $u_i>0$ for at least one
$i\in L$.  Equivalently its system carries a strictly positive coefficient on
some \emph{spoke} $H\cup\{i\}$, $i\in L$; in particular $b>0$.
\end{lemma}

\begin{proof}
Choose $z\in\core(v)$ with $z(H)=v(H)$ and put $\Delta=z-\phi$.  Each supported
$T$ is tight and feasible, so $\Delta(T)\ge0$; each spoke satisfies
$\Delta(H\cup\{i\})\ge-s_\phi(H\cup\{i\})$; and $\Delta(H)=-s$.  Evaluating
$\Lambda=b\one_N$ at $\Delta$ and using $\Delta(N)=0$ gives (PB).  If $u=0$,
(PB) reads $\delta s\le0$, contradicting $\delta>0<s$.  Finally $u_i>0$ forces
$b>c_i\ge0$.
\end{proof}

\begin{corollary}\label{sel:cor:spoke}
Let a permutational type have exceptional set $E$ and no positive set of the
form $E\cup\{i\}$.  Then it is the type of no normal certificate at a slack
boundary, hence of no packet and of no mixture.  Every type with $r=0$ has this
property: if $r=0$ and $i\notin E$ then the load at $i$ vanishes while every
set containing $i$ has nonnegative coefficient, so every positive set lies
inside $E$.
\end{corollary}

The second lemma is an uncrossing inequality for tight coalitions.  It was
obtained independently in the companion five- and six-player computations, and
is stated here once.

\begin{lemma}[crossing budget]\label{sel:lem:crossing}
Let $v$ be exact and singleton-normalized and let $\phi\in\core(v)$.  If $S_1$
and $S_2$ are $\phi$-tight then
\[
 s_\phi(S_1\cap S_2)\ \le\ \phi\bigl(N\setminus(S_1\cup S_2)\bigr).
\]
\end{lemma}

\begin{proof}
Put $C=S_1\cap S_2$, $U=S_1\cup S_2$ and $R=N\setminus U$.  Exactness supplies
$z\in\core(v)$ with $z(C)=v(C)$, whence
\[
 \phi(S_1)+\phi(S_2)=v(S_1)+v(S_2)\le z(S_1)+z(S_2)=z(U)+z(C)
   =\phi(N)-z(R)+v(C).
\]
Since $\phi(S_1)+\phi(S_2)=\phi(U)+\phi(C)$ and $\phi(U)=\phi(N)-\phi(R)$, this
is $s_\phi(C)\le\phi(R)-z(R)\le\phi(R)$, the last step by $z\ge0$.
\end{proof}

Two specializations are used.  If $S_1\cup S_2=N$ then $S_1\cap S_2$ is tight.
If $S_1\cap S_2=\{j\}$ then $\phi_j\le\phi(N\setminus(S_1\cup S_2))$; taking
$j\in H$ and $S_1\cup S_2=N\setminus\{i\}$ with $i\in L$ contradicts the strict
boundary.  The four-player arguments of the companion note are the case
$S_1\cup S_2=N$.

The third lemma is the exact dual of singleton normalization and sharpens
Lemma~\ref{lem:overload} at a singleton face.

\begin{lemma}[zeroing obstruction]\label{sel:lem:zero}
Let $v$ be exact and singleton-normalized, $\phi\in\core(v)$, and let
$\lambda\ge0$ be weights on proper $\phi$-tight coalitions whose loads satisfy
$c_k\le1$ for every $k\ne j$.  Then $\sum_kc_k\phi_k\le\phi(N)$; equivalently,
with $c_j=1+\theta$,
\[
 \theta\,\phi_j\ \le\ \sum_{k\ne j}(1-c_k)\phi_k .
\]
\end{lemma}

\begin{proof}
Exactness at $\{j\}$ gives $z\in\core(v)$ with $z_j=v(\{j\})=0$.  Then
$\sum_T\lambda_T\phi(T)=\sum_T\lambda_Tv(T)\le\sum_T\lambda_Tz(T)
=\sum_kc_kz_k=\sum_{k\ne j}c_kz_k\le\sum_{k\ne j}z_k=\phi(N)$,
and the left side is $\sum_kc_k\phi_k$.
\end{proof}

\begin{corollary}\label{sel:cor:zero}
If $j$ is a richest player and $\sum_T\lambda_T|T|\ge|N|$, then $\phi_k=\phi_j$
for every $k$ with $c_k<1$.  In particular the configuration is impossible as
soon as some $k$ with $c_k<1$ has $\phi_k<\phi_j$.
\end{corollary}

By linear-programming duality, Lemma~\ref{sel:lem:zero} is not merely necessary:
a core allocation vanishing at $j$ fails to exist precisely when some such
$\lambda$ has $\sum_kc_k\phi_k>\phi(N)$.  It is the only obstruction that is
ever active in the classification below.

\subsection{Full-trace packets are stars, and stars are facets}

\begin{definition}[star system]\label{sel:def:star}
For $\emptyset\ne H\subsetneq N$, $L=N\setminus H$ and $\emptyset\ne
Q\subsetneq L$, put
\[
 \mathrm{St}(H,Q):\qquad
 \one_{H\cup Q}+\sum_{i\in L\setminus Q}\one_{H\cup\{i\}}
   -(|L|-|Q|)\,\one_H=\one_N .
\]
\end{definition}

\begin{lemma}[star normal form]\label{sel:lem:star}
Let the boundary be slack.
\begin{enumerate}[nosep,label=(\roman*)]
\item If $h=1$ then $\mathcal P_H$ has the single vertex $\beta_H=1$, so every
      packet has exactly one lift $T$; $T$ is proper and $T\ne H$, hence
      $T=H\cup Q$ with $\emptyset\ne Q\subsetneq L$, and the system is
      $\mathrm{St}(H,Q)$, with $b=1$ and $\delta=m-|Q|$.
\item For arbitrary $h$, every \emph{full-trace} packet has one lift
      $T\supseteq H$, again proper and $\ne H$, so its system is
      $\mathrm{St}(H,T\cap L)$.
\end{enumerate}
In either case every positive coefficient equals one and $b=1$.
\end{lemma}

\begin{proposition}[stars are facets]\label{sel:prop:starfacet}
$\mathrm{St}(H,Q)$ is a minimal semi-balanced system with exceptional set $H$,
and it is indecomposable.  Its permutational type depends only on
$(h,q)=(|H|,|Q|)$, and the admissible pairs are $h\ge1$, $q\ge1$,
$h+q\le n-1$, so there are $\binom{n-1}{2}$ star types on $n$ players.
\end{proposition}

\begin{proof}
Write $\mathcal S=\{H\}\cup\{H\cup\{i\}:i\in L\setminus Q\}\cup\{H\cup Q\}$,
$m=|L|$, $q=|Q|$, and note $1\le q\le m-1$, so $L\setminus Q\ne\emptyset$ and
$H\cup Q\ne N$.  For $k\in H$ every member of $\mathcal S$ contains $k$; for
$k\in Q$ the only member containing $k$ is $H\cup Q$; for $k\in L\setminus Q$
the only member containing $k$ is $H\cup\{k\}$.  The displayed coefficients are
all nonzero with exactly one negative and give load one everywhere, so the
system is semi-balanced with $r=1$.

\emph{Minimality.}  Let $\emptyset\ne\mathcal S'\subsetneq\mathcal S$ carry
nonzero coefficients $\mu$, at most one negative, with constant load $r'$.
Reading the load at a player of $Q$ and at a player of $L\setminus Q$ shows
that $\mu_{H\cup Q}=r'$ if $H\cup Q\in\mathcal S'$ and $r'=0$ otherwise, and
likewise $\mu_{H\cup\{i\}}=r'$ if $H\cup\{i\}\in\mathcal S'$ and $r'=0$
otherwise.  If $r'=0$ then no set other than $H$ can lie in $\mathcal S'$
without a vanishing coefficient, so $\mathcal S'=\{H\}$ and the load at $H$
gives $\mu_H=0$, excluded.  If $r'\ne0$ then $H\cup Q$ and every
$H\cup\{i\}$ lies in $\mathcal S'$ with coefficient $r'$, so
$\mathcal S'=\mathcal S\setminus\{H\}$ by properness, and the load at a player
of $H$ gives $(m-q+1)r'=r'$, i.e.\ $r'=0$ since $q<m$.  Contradiction.

\emph{Indecomposability.}  Suppose $E$ is a proper nonempty set with
$E\notin\mathcal S$ and there are $\sigma\ge0$ on $\mathcal S$ and $t>0$ with
$\sum_{S\in\mathcal S}\sigma_S\one_S-t\one_E=r\one_N$.  Write
$x=\sigma_H$, $y_i=\sigma_{H\cup\{i\}}$, $z=\sigma_{H\cup Q}$,
$\Sigma=x+\sum_iy_i+z$ and $\varepsilon_k=[k\in E]$.  The three families of
load equations read
\[
 \Sigma-t\varepsilon_k=r\ (k\in H),\qquad
 z-t\varepsilon_k=r\ (k\in Q),\qquad
 y_k-t\varepsilon_k=r\ (k\in L\setminus Q).
\]
The first shows $\varepsilon$ is constant on $H$, and the second that it
is constant on $Q$; so $E\cap H\in\{\emptyset,H\}$ and
$E\cap Q\in\{\emptyset,Q\}$.

If $E\cap H=\emptyset$ then $r=\Sigma$.  For $k\in Q$ we get
$z-t\varepsilon_k=\Sigma\ge z$, so $\varepsilon_k=0$ and $z=\Sigma$, whence
$x=0$ and every $y_i=0$.  Then for $k\in L\setminus Q$,
$-t\varepsilon_k=r=z$; positivity of $t$ forces $z=0$ and
$\varepsilon_k=0$, so $E=\emptyset$, a contradiction.

If $H\subseteq E$ then $r=\Sigma-t$.  Suppose first $Q\subseteq E$.  Then
$z=\Sigma$, so $x=0$ and every $y_i=0$; for $k\in L\setminus Q$,
$-t\varepsilon_k=z-t$.  If some $\varepsilon_k=1$ then $z=0$, all $\sigma$
vanish and $-t\one_E=r\one_N$ forces $E\in\{\emptyset,N\}$, excluded; so every
$\varepsilon_k=0$ and $E=H\cup Q\in\mathcal S$, excluded.  Suppose instead
$Q\cap E=\emptyset$.  Then $z=\Sigma-t$, i.e.\ $x+\sum_iy_i=t>0$.  If
$\varepsilon_k=1$ for some $k\in L\setminus Q$ then $y_k=\Sigma$, so all other
weights vanish, hence $\varepsilon_i=0$ and $y_i=r=0$ for $i\ne k$, and
$E=H\cup\{k\}\in\mathcal S$, excluded.  If $\varepsilon_k=0$ for every
$k\in L\setminus Q$ then $E=H\in\mathcal S$, excluded.

The remaining assertions are clear: relabelling identifies any two stars with
the same $(h,q)$, and $q\le m-1=n-h-1$.
\end{proof}

\subsection{The five-player classification}

At $n=5$ Theorem~\ref{thm:phase} admits only $(m,h)=(4,1)$ and $(3,2)$.  For
those two shapes one may enumerate, with no game and no payoffs, every system
(P) with $\delta>0$ that Theorem~\ref{thm:packet} can produce: the traces range
over the minimal balanced collections on $H$ and each lift over every proper set
with the prescribed trace that is not $H$ and not a singleton inside $H$.  There
are three types at $(4,1)$, all full-trace, and eight at $(3,2)$, of which two
are full-trace and six are split.

\begin{theorem}[five-player selection]\label{sel:thm:n5}
Let $v$ be a singleton-normalized exact five-player game with a slack strict
boundary $L|H$ at $\phi=\Lmin(v)$.  Then a system (P) is a facet of the cone of
exact games---equivalently an indecomposable minimal semi-balanced system,
equivalently a member of the catalogue---if and only if its packet is
full-trace, in which case it is the star $\mathrm{St}(H,T\cap L)$.  The
facet-defining types that occur are exactly
\[
 \begin{array}{llll}
  (h,q)=(1,1)\to\text{row }16, &(1,2)\to\text{row }14, &(1,3)\to\text{row }11,\\
  (h,q)=(2,1)\to\text{row }13, &(2,2)\to\text{row }10, &
 \end{array}
\]
and all five are realized by integral exact games.  Systems (P) that are not
facets also occur; see Example~\ref{sel:ex:split}.
\end{theorem}

The eleven excluded rows fall into three groups of visibly different depth.

\begin{itemize}[nosep]
\item Rows $1$--$8$ and row $12$ are excluded by Corollary~\ref{sel:cor:spoke}:
      none of them carries a positive spoke over its exceptional set.  Rows
      $1$--$8$ have $r=0$; row $12$ is $-a+abc+ade$, whose exceptional set is
      $\{a\}$ and which has no positive set of size two containing $a$.  This
      exclusion is robust: those types are the systems of no normal
      certificate, hence of no mixture either.
\item Row $9$ is excluded by Theorem~\ref{thm:phase} alone.  It is the star with
      $(h,q)=(3,1)$, and as such is a perfectly legitimate packet shape; it
      needs $(m,h)=(2,3)$, which $m_{\min}(3)=3$ forbids on five players.
\item Row $15$ is excluded by Lemma~\ref{sel:lem:star}(i): an $h=1$ packet has
      $b=1$ and unit positive coefficients, whereas row $15$ has $r=2$ and a
      coefficient $2$.
\end{itemize}

\begin{remark}[what the selection, rather than the geometry, excludes]
\label{sel:rem:mixture}
The exclusions of rows $12$ and $15$ are both statements about systems (P), and
both are correct, but they are not equally deep.  Row $15$ \emph{is} the system
of a normal certificate: at $h=1$ every $\phi$-tight coalition has trace $H$, so
a cover may support several lifts of the one trace, and
$\lambda_{acd}=\lambda_{ace}=\lambda_{ade}=\tfrac13$ has rich load one and
poorer loads $(0,\tfrac23,\tfrac23,\tfrac23)$, all below one, giving
$b=\tfrac23$, $u_b=\tfrac23$, $\delta=1$ and exactly the row-$15$ system.  What
the one-packet theorem then does is replace that cover by one of its three
lifts, reporting a row-$14$ star.  Row $12$ admits no such certificate, by
Corollary~\ref{sel:cor:spoke}.  In short: row $12$ is excluded by the geometry,
row $15$ only by the selection.
\end{remark}

\subsection{Split packets, and a threshold in the number of poorer players}

The full-trace conclusion does \emph{not} extend to split packets, even on five
players.

\begin{example}[a split packet with $\delta>0$]\label{sel:ex:split}
Apply Lemma~\ref{lem:envelope} to the five common-total rows
\[
 (0,1,0,3,3),\quad(1,0,0,3,3),\quad(1,1,1,2,2),\quad
 (2,2,0,0,3),\quad(2,2,0,3,0),
\]
equivalently take the core polytope
$\{x\ge0:\ x(N)=7,\ x(abe)\ge4,\ x(abd)\ge4,\ x(ade)\ge5,\ x(bde)\ge5\}$.
The game is exact and singleton-normalized, $\Lmin(v)=(1,1,1,2,2)$, the strict
boundary is $L=\{a,b,c\}$, $H=\{d,e\}$, and $v(H)=3<4=\phi(H)$, so $s=1$.  The
proper $\phi$-tight coalitions are $abd$, $abe$, $ade$, $bde$, $abde$, whose
traces on $H$ are $\{d\}$, $\{e\}$ and $\{d,e\}$: \emph{both} proper traces
occur.  The split vertex lifted by $T_1=abd$ and $T_2=abe$ has
$r=(2,2,0)$, $b=2$, $u_c=2$ and $\delta=1>0$, with system
\[
 \one_{abd}+\one_{abe}+2\,\one_{cde}-\one_{de}=2\,\one_N .
\]
This system is minimal semi-balanced but decomposable, by $E=\{d\}$:
$\one_{abd}+\one_{cde}-\one_d=\one_N$.  The same boundary also carries the
full-trace packets of types row~$13$ and row~$10$.
\end{example}

Split packets are nevertheless invisible to the catalogue as long as the poorer
side is small.  The next statement is a finite check on set systems: it
enumerates every abstract system (P) with $\delta>0$ for the stated $(m,h)$ and
classifies each permutational type, and therefore applies at every boundary of
every exact game.

By Proposition~\ref{prop:m2} an $h=2$ boundary needs $m\ge3$, so the statement
below starts where Lemma~\ref{lem:rankthree} stops.

\begin{theorem}[split packets and the poorer-side threshold]\label{sel:thm:mthr}
Let $h=2$.  If $m\le4$ then every split packet system is minimal but
decomposable; consequently every facet-defining packet at an $h=2$ boundary
with $m\le4$ is a full-trace star.  This fails at $m=5$.  Precisely:
\begin{center}
\begin{tabular}{lrrrr}
\toprule
 & packets & full-trace types & split types & split facets\\
\midrule
$m=3$ ($n=5$) &  42 & 2 &  6 & 0\\
$m=4$ ($n=6$) & 216 & 3 & 13 & 0\\
$m=5$ ($n=7$) & 950 & 4 & 24 & 2\\
\bottomrule
\end{tabular}
\end{center}
The two exceptions at $m=5$ are
$af-5\,fg+bfg+cfg+2\,dfg+2\,efg+abcg$ and
$-5\,fg+abf+acg+bfg+cfg+2\,dfg+2\,efg$.
Equivalently, no catalogued type with a doubleton exceptional set has
split-packet shape when $m\le4$: on five players the doubleton rows are
$1,10,13$, of which $10$ and $13$ are stars and $1$ has no spoke; on six players
they are $1,25,28,29,33,34,39$, of which $28,33,39$ are stars and none is split.
\end{theorem}

The same finite check explains the analogous phenomenon one dimension up.

\begin{theorem}[which trace collections carry facets at $h=3$]\label{sel:thm:h3}
Let $(m,h)=(3,3)$, so $n=6$.  Among the $811$ abstract packets with
$\delta>0$, sorted by the minimal balanced collection their traces form:
\begin{center}
\begin{tabular}{lrrl}
\toprule
collection on $H$ & types & facets & rows\\
\midrule
$\{H\}$ (full trace)       &  2 & 2 & $32,\ 27$\\
three pairs (the 2-cover)   & 31 & 5 & $7,\ 24,\ 51,\ 52,\ 53$\\
a singleton and a pair      & 10 & 0 & ---\\
three singletons            & 12 & 0 & ---\\
\bottomrule
\end{tabular}
\end{center}
Only the full trace and the two-cover carry facet-defining systems.
\end{theorem}

Theorems~\ref{sel:thm:mthr} and~\ref{sel:thm:h3} share one mechanism: the
partition-type minimal balanced collections generate decomposable systems, so
they can certify a failure but never a facet of the exact cone.  The
five-player classification is the case where that mechanism is complete.

\subsection{What is proved, what is checked, and what is verified}

\begin{center}
\begin{tabular}{lp{0.62\textwidth}}
\toprule
\textbf{status} & \textbf{statement}\\
\midrule
proved &
 Lemma~\ref{sel:lem:spoke} and Corollary~\ref{sel:cor:spoke};
 Lemma~\ref{sel:lem:crossing}; Lemma~\ref{sel:lem:zero} and
 Corollary~\ref{sel:cor:zero}; Lemma~\ref{sel:lem:star};
 Proposition~\ref{sel:prop:starfacet}; the exclusion of rows $1$--$8$, $9$,
 $12$, $15$ on five players; Remark~\ref{sel:rem:mixture}.\\
finite check &
 the list of eleven abstract packet shapes on five players;
 Theorem~\ref{sel:thm:mthr}; Theorem~\ref{sel:thm:h3}.  Each enumerates set
 systems only, so each holds at every boundary of every exact game.\\
verified &
 Example~\ref{sel:ex:split} and the five realizing games of
 Theorem~\ref{sel:thm:n5}, by exact rational recomputation of the envelope, of
 $\Lmin$, of the tight family and of every packet.  That exactly one of the six
 split shapes on five players is realizable rests on a search over tight
 families of at most $|N|+2$ coalitions, which is exhaustive when the poorer
 side is a single payoff stratum and a broad search otherwise.\\
\bottomrule
\end{tabular}
\end{center}

Two caveats are owed.  First, the claim that the split shape of
Example~\ref{sel:ex:split} is the \emph{only} realizable one is the one
statement above that is not a theorem; the five excluded shapes survive the
tangent-cone conditions of Lemma~\ref{lem:cover} and are killed only by
singleton normalization, i.e.\ by Lemma~\ref{sel:lem:zero}, and we have no
human proof of that.  Second, earlier random harvests reported no split packets
at all on five players; Example~\ref{sel:ex:split} shows that this was a
sampling artefact, and the same caution applies to any statement of the form
``collection $\mathcal R$ is never selected'' that rests on random exact games
rather than on a finite check.

\section{What the theorem does and does not establish}

The normal-cover and one-packet theorems are dimension-free: every failure is
reduced to a single minimal balanced rich-trace packet and its poorer-side
completion budget.  The phase theorem then classifies every possible pair of
boundary-side cardinalities, and Section~\ref{sec:selection} classifies the
selected systems completely on five players and by finite check for
$(m,h)\le(4,2)$ and $(3,3)$.  What remains unclassified is the general
dimension: distinct irreducible packet classes occur (full-trace stars, the
two-cover class of Theorem~\ref{sel:thm:h3}, and from $m=5$ on, facet-defining
split classes), so a complete catalogue is multi-class, not a theorem that
every failure contains one fixed finite configuration.

One calibration is owed to the reader.  The linear-size bound in
Corollary~\ref{cor:certificate}(i) is not itself new: it is the standard fact
that a basic solution of a system of equalities has support at most the
number of rows, applied to (2), whose rows are indexed by $H$.  Sparse
supports of this kind are routine in the balancedness literature surrounding
the Kohlberg criterion.  What is claimed here is not the bound but which
single system the Lmin geometry selects, and on which prescribed suffix it
lives.

The theorem's novelty remains conditional on specialist review.  Exact
partition, its Lorenz consequence, property I, semi-balanced inequalities,
the semi-balanced/minimal-balanced correspondence, minimal balanced
collections, and lower-envelope exactification are not claimed.  The
candidate new content is the simultaneous Lmin normal cover, the richer-side
one-packet selection with its linear-size witness and slack-propagation
corollary, the sharp staircase (1), the three dimension-free
indecomposable realizing families, and the selection results of
Section~\ref{sec:selection}: the positive-load, crossing-budget and zeroing
lemmas, the star normal form with its facet proposition, the five-player
selection theorem, and the poorer-side threshold for split facets.

\section*{Reproducibility note}

The proof above is human and enumerates neither games nor payoff-stratum
patterns.  Exact
rational programs in the companion repository independently reconstruct the
smallest family members, all coalition values, and staged Lmin solutions, and
classify the packet systems of Remark~\ref{rem:catalogue} against the
published catalogue of semi-balanced types.
They are bug audits, not proof evidence.

\begin{thebibliography}{99}

\bibitem[Arin and I\~narra(2001)]{ArinInarra2001}
Arin J, I\~narra E (2001) Egalitarian solutions in the core.
International Journal of Game Theory 30:187--193.

\bibitem[Arin et al.(2003)]{ArinKuipersVermeulen2003}
Arin J, Kuipers J, Vermeulen D (2003) Some characterizations of egalitarian
solutions on classes of TU-games. Mathematical Social Sciences 46:327--345.
\url{https://doi.org/10.1016/S0165-4896(03)00048-6}

\bibitem[Biswas et al.(1999)]{Biswas1999}
Biswas AK, Parthasarathy T, Potters JAM, Voorneveld M (1999) Large cores and
exactness. Games and Economic Behavior 28:1--12.
\url{https://doi.org/10.1006/game.1998.0683}

\bibitem[Cs\'oka et al.(2011)]{CsokaHeringsKoczy2011}
Cs\'oka P, Herings PJJ, K\'oczy L\'A (2011) Balancedness conditions for exact
games. Mathematical Methods of Operations Research 74:41--52.

\bibitem[Dietzenbacher(2021)]{Dietzenbacher2021}
Dietzenbacher B (2021) Monotonicity and egalitarianism.
Games and Economic Behavior 127:194--205.
\url{https://doi.org/10.1016/j.geb.2021.03.006}

\bibitem[Llerena and Mauri(2017)]{LlerenaMauri2017}
Llerena F, Mauri L (2017) On the existence of the Dutta--Ray's egalitarian
solution. Mathematical Social Sciences 89:92--99.
\url{https://doi.org/10.1016/j.mathsocsci.2017.07.002}

\bibitem[Laplace Mermoud et al.(2025)]{LaplaceMermoudEtAl2025}
Laplace Mermoud L, Grabisch M, Sudh\"olter P (2025) Minimal balanced
collections and their applications to core stability and other topics of
game theory. Working paper.

\bibitem[Lohmann et al.(2012)]{LohmannBormHerings2012}
Lohmann E, Borm P, Herings PJJ (2012) Minimal exact balancedness.
Mathematical Social Sciences 64:127--135.

\bibitem[Mauri(2018)]{Mauri2018}
Mauri Masdeu L (2018) Essays on egalitarianism-based solution concepts for
cooperative TU-games. Doctoral thesis, Universitat Rovira i Virgili.

\bibitem[Studen\'y and Kratochv\'il(2022)]{StudenyKratochvil2022}
Studen\'y M, Kratochv\'il V (2022) Facets of the cone of exact games.
Mathematical Methods of Operations Research 95:35--80.
\url{https://doi.org/10.1007/s00186-022-00770-4}.
Electronic catalogue of indecomposable min-semi-balanced systems:
\url{http://gogo.utia.cas.cz/indecomposable-min-semi-balanced-catalogue/}

\bibitem[S\'anchez-Soriano et al.(2010)]{SanchezSorianoBranzeiLlorcaTijs2010}
S\'anchez-Soriano J, Br\^anzei R, Llorca N, Tijs SH (2010) A technical note
on Lorenz dominance in cooperative games. CentER Discussion Paper 2010--101.
\url{https://pure.uvt.nl/ws/files/1269376/2010-101.pdf}

\end{thebibliography}

\end{document}
